\documentclass[12pt]{article}
\usepackage[T1]{fontenc}
\usepackage[utf8]{inputenc}
\usepackage{lmodern}
\usepackage{textcomp}
\usepackage{enumitem}
\usepackage{geometry}
\geometry{margin=1in}
\begin{document}
\begin{center}
Answer Key - Chemistry - Graduate Level Assessment 2 \
\small (Answers are ordered exactly as in the question paper.)
\end{center}

\section*{Multiple Choice}
\begin{enumerate}[label=\arabic*.]
\item \textbf{Answer:} D
\\ \textit{Options:} A) Eu; B) Gd; C) Am; D) Cu
\\ \textit{Note:} Copper (Cu) has an electron configuration of [Ar] 3 $d^10$ 4 $s^1$, indicating its 4 f and 5 f orbitals are fully filled, unlike the lanthanides and actinides that possess partially filled 4 f or 5 f orbitals.
\item \textbf{Answer:} A
\\ \textit{Options:} A) 10 ml; B) 30 ml; C) 55 ml; D) 40 ml; E) 25 ml
\\ \textit{Note:} The correct answer is 10 ml because the stoichiometry of the neutralization reaction indicates that 2 moles of H$_{3}$PO$_{4}$ are required to neutralize 3 moles of Ca(OH)2, leading to the calculation that 0.025 M H$_{3}$PO$_{4}$ needs to be 10 ml to completely react with the 25 ml of 0.030 M Ca(OH)2.
\item \textbf{Answer:} A
\\ \textit{Options:} A) 1.6 \ensuremath{\times} 10^-3 F; B) 1.5 \ensuremath{\times} 10^-3 F; C) 1.4 \ensuremath{\times} 10^-3 F; D) 0.8 \ensuremath{\times} 10^-3 F; E) 0.9 \ensuremath{\times} 10^-3 F
\\ \textit{Note:} The correct answer is 1.6 × 10^-3 F because the formality is calculated by determining the moles of acetic acid in one liter of the wine solution, which involves using the weight percentage of acetic acid, the density of the wine, and the formula weight of acetic acid to find the concentration in terms of moles per liter.
\item \textbf{Answer:} E
\\ \textit{Options:} A) .1002 : .3002 : .4002; B) .3582 : .4582 : .0582; C) .2582 : .4998 : .3092; D) .1582 : .6008 : .2092; E) .1582 : .4998 : .1092
\\ \textit{Note:} The correct answer reflects Graham's law of effusion, which states that the rate of diffusion of a gas is inversely proportional to the square root of its molar mass, resulting in the observed ratios for Ar, He, and Kr based on their respective molar masses.
\item \textbf{Answer:} B
\\ \textit{Options:} A) He\_2^\ensuremath{\oplus} is the most stable among He, He\_2, and He\_2^\ensuremath{\oplus}; B) He\_2^\ensuremath{\oplus} is more stable than He\_2 but less stable than He; C) He, He\_2, and He\_2^\ensuremath{\oplus} have the same stability; D) He\_2^\ensuremath{\oplus} is the least stable among He, He\_2, and He\_2^\ensuremath{\oplus}; E) He and He\_2 have the same stability, both more stable than He\_2^\ensuremath{\oplus}
\\ \textit{Note:} He\_2^\ensuremath{\oplus} is more stable than He\_2 because the removal of an electron from He\_2 reduces electron-electron repulsion and results in a greater bond order, while He is more stable than both due to its complete electron configuration and higher bond energy.
\end{enumerate}

\section*{True / False}
\begin{enumerate}[label=\arabic*.]
\setcounter{enumi}{5}
\item \textbf{Answer:} True
\\ \textit{Options:} A) True; B) False
\\ \textit{Note:} The ionization energy of singly ionized helium is accurately calculated as 48.359 eV using the Bohr model due to its dependence on the effective nuclear charge and the formula for energy levels in hydrogen-like atoms, which reflects the increased attraction of the nucleus to the electron in helium.
\item \textbf{Answer:} True
\\ \textit{Options:} A) True; B) False
\\ \textit{Note:} At$_{12}$.19 K, an ideal gas can exist at 1 atm pressure and a concentration of 1 mole liter^-1 because the ideal gas law (PV=nRT) holds true under these conditions, as long as the gas is not at a temperature or pressure where deviations from ideal behavior occur.
\item \textbf{Answer:} False
\\ \textit{Options:} A) True; B) False
\\ \textit{Note:} Zn(NO 3)2 is actually more soluble in water than NiSO 3 because nitrate ions (NO 3-) are highly soluble in water, whereas the sulfate ion (SO 3) tends to form less soluble compounds with certain cations.
\item \textbf{Answer:} False
\\ \textit{Options:} A) True; B) False
\\ \textit{Note:} The missing mass in ^19\_9 F, which corresponds to the mass defect due to binding energy, is not significantly greater than the mass of an electron, as the mass defect is typically much smaller than 0.199 amu for light nuclei like fluorine.
\item \textbf{Answer:} False
\\ \textit{Options:} A) True; B) False
\\ \textit{Note:} The statement is false because the molality is calculated using the number of moles of solute (H$_{2}$SO$_{4}$) divided by the mass of the solvent (water) in kilograms, and with 49 g of H$_{2}$SO$_{4}$ corresponding to approximately 0.49 moles, the molality is actually 1.96 moles/kg, not 1.0 moles/kg.
\end{enumerate}

\section*{Long Form Response}
\begin{enumerate}[label=\arabic*.]
\setcounter{enumi}{10}
\item \textbf{Answer:} To dilute 5 gallons of a 90\% ammonia solution to a 60\% concentration, we need to calculate the volume of water required. First, we determine the amount of ammonia in the initial solution: 5 gallons \ensuremath{\times} 0.90 = 4.5 gallons of ammonia. Let x be the volume of water added. After dilution, the total volume will be 5 + x gallons, and we want the ammonia concentration to be 60\%. The equation can be set up as follows: 4.5 gallons of ammonia / (5 + x) gallons = 0.60. Solving for x, we find that 4.5 = 0.60(5 + x), leading to 4.5 = 3 + 0.60 x, which simplifies to 1.5 = 0.60 x. Thus, x = 2.5 gallons. Therefore, 2.5 gallons of water are needed to achieve the desired 60\% ammonia concentration.
\\ \textit{Note:} correct because it accurately applies the concept of dilution by calculating the initial amount of ammonia, setting up a proportion to find the final concentration, and solving for the volume of water needed to achieve the desired 60\% concentration.
\item \textbf{Answer:} The expansion of one cubic foot of air when heated from 0^\ensuremath{\circC} to 546,000^\ensuremath{\circC} can be analyzed through the principles of thermodynamics, particularly the Ideal Gas Law, which states that PV = nRT. As temperature increases significantly, the volume of the gas expands exponentially, assuming pressure remains constant. In this scenario, the resulting volume of the air expands to approximately 3000 cubic feet due to the drastic increase in temperature, causing the gas molecules to gain kinetic energy and occupy greater space. This dramatic change illustrates the profound impact of thermal energy on gas behavior, emphasizing the relationship between temperature and volume as described by Charles's Law. The calculations underlying this expansion demonstrate not only the theoretical principles but also the practical implications of extreme thermal conditions on gases.
\\ \textit{Note:} correct because it applies the Ideal Gas Law, illustrating that the significant increase in temperature from 0°C to 546,000°C leads to an exponential increase in volume, assuming constant pressure, due to the heightened kinetic energy of the gas molecules.
\end{enumerate}

\vfill
\noindent\textit{End of Answer Key}
\end{document}